%%%%%%%%%%%%%%%%%%%%%%%%%%%%%%%%%%%%%%%%%%%%%%%%%%%%%%%%%%%%%%%%%%%%%%%%%%
%%%%%%%%%%%%%%%%%%%%%%%%%%     Cilindro N Cone      %%%%%%%%%%%%%%%%%%%%%%%
%%%%%%%%%%%%%%%%%%%%%%%%%%%%%%%%%%%%%%%%%%%%%%%%%%%%%%%%%%%%%%%%%%%%%%%%%%

% Corte, num plano coordenado, de um cilindro de raio a e de um cone reto
% (abertura de 45°): as retas h = -a e h = a, as retas v = h e v = -h, os
% pontos de interseção e parte da região dentro do cilindro e fora do cone
% (|v| <= |h| <= a, os dois triângulos deitados) sombreada.
% Uso: \CilindroNcone{eixo horizontal}{eixo vertical}{a}{q}
% Ex.: \CilindroNcone{x}{z}{2}{1}  % triângulo do 1º quadrante, no plano xz
% q = 1, 2, 3, 4 : parte da região num quadrante
% q = t, b, d, e : metade de cima (top), de baixo (bottom), direita ou esquerda
%                  (l e r também valem, como no \Aro)
% q = o          : a região toda

\makeatletter

% Comando auxiliar, com a região definida explicitamente
% #1, #2: eixos (horizontal e vertical)   #3: raio do cilindro (a)
% #4, #5, #6, #7: retângulo de recorte do sombreado, em múltiplos de 1.3a:
%   h de #4 a #5 e v de #6 a #7 (cada um -1, 0 ou 1)
\newcommand{\CilindroNconeAux}[7]{%
\begin{tikzpicture}[scale=2]
  \begin{axis}[
      xlabel={$#1$}, ylabel={$#2$},
      axis equal image,             % mantém escala para os 45° ficarem corretos
      xmin=-1.3*#3, xmax=1.3*#3, ymin=-1.3*#3, ymax=1.3*#3,
      clip=false,                   % permite que rótulos ultrapassem a área dos eixos
      xtick={-#3,#3}, ytick={-#3,#3},
      x tick label style={anchor=north east}, % rótulos abaixo e à esquerda do eixo horizontal
      y tick label style={anchor=south east}, % rótulos acima e à esquerda do eixo vertical
      every axis x label/.style={at={(ticklabel* cs:0.97)}, anchor=north},
      every axis y label/.style={at={(ticklabel* cs:0.97)}, anchor=east},
    ]

    % --- Linhas tracejadas ---
    \addplot[Tracejado] coordinates {(#3,0) (#3,#3)};
    \addplot[Tracejado] coordinates {(0,#3) (#3,#3)};

    % --- Curvas ---
    \addplot[thick, blue, domain=-1.2*#3:1.2*#3] ({#3}, {x}) node[pos=0.85, right, blue] {$#1=#3$};   % Reta h = a
    \addplot[thick, blue, domain=-1.2*#3:1.2*#3] ({-#3}, {x}) node[pos=0.85, left, blue] {$#1=-#3$};  % Reta h = -a
    \addplot[thick, red, domain=-1.2*#3:1.2*#3, samples=2] {x} node[pos=1, above, red] {$#2=#1$};     % Reta v = h
    \addplot[thick, red, domain=-1.2*#3:1.2*#3, samples=2] {-x} node[pos=1, below, red] {$#2=-#1$};   % Reta v = -h

    % --- Região sombreada, recortada pela parte escolhida (q) ---
    \begin{scope}
    \clip (axis cs:{#4*1.3*#3},{#6*1.3*#3}) rectangle (axis cs:{#5*1.3*#3},{#7*1.3*#3});
    \path[fill=gray!30, opacity=0.6]                                % Triângulo da direita
      (axis cs:0,0) -- (axis cs:#3,#3) -- (axis cs:#3,-#3) -- cycle;
    \path[fill=gray!30, opacity=0.6]                                % Triângulo da esquerda
      (axis cs:0,0) -- (axis cs:-#3,#3) -- (axis cs:-#3,-#3) -- cycle;
    \end{scope}

    % --- Pontos de interseção ---
    \addplot[only marks, black!60, mark=*] coordinates {(#3,#3) (-#3,#3) (#3,-#3) (-#3,-#3)};

  \end{axis}
\end{tikzpicture}%
}

% Tabela de q -> retângulo de recorte. Cada entrada recebe #1, #2 (eixos) e #3 (a).
\@namedef{CilindroNcone@1}#1#2#3{\CilindroNconeAux{#1}{#2}{#3}{0}{1}{0}{1}}
\@namedef{CilindroNcone@2}#1#2#3{\CilindroNconeAux{#1}{#2}{#3}{-1}{0}{0}{1}}
\@namedef{CilindroNcone@3}#1#2#3{\CilindroNconeAux{#1}{#2}{#3}{-1}{0}{-1}{0}}
\@namedef{CilindroNcone@4}#1#2#3{\CilindroNconeAux{#1}{#2}{#3}{0}{1}{-1}{0}}
\@namedef{CilindroNcone@t}#1#2#3{\CilindroNconeAux{#1}{#2}{#3}{-1}{1}{0}{1}}
\@namedef{CilindroNcone@b}#1#2#3{\CilindroNconeAux{#1}{#2}{#3}{-1}{1}{-1}{0}}
\@namedef{CilindroNcone@d}#1#2#3{\CilindroNconeAux{#1}{#2}{#3}{0}{1}{-1}{1}}
\@namedef{CilindroNcone@e}#1#2#3{\CilindroNconeAux{#1}{#2}{#3}{-1}{0}{-1}{1}}
\@namedef{CilindroNcone@r}#1#2#3{\CilindroNconeAux{#1}{#2}{#3}{0}{1}{-1}{1}}
\@namedef{CilindroNcone@l}#1#2#3{\CilindroNconeAux{#1}{#2}{#3}{-1}{0}{-1}{1}}
\@namedef{CilindroNcone@o}#1#2#3{\CilindroNconeAux{#1}{#2}{#3}{-1}{1}{-1}{1}}

\newcommand{\CilindroNcone}[4]{%
  \@ifundefined{CilindroNcone@#4}%
    {\PackageError{CilindroNcone}{q = '#4' invalido}{Use 1, 2, 3, 4, t, b, d, e (ou l, r) ou o.}}%
    {\@nameuse{CilindroNcone@#4}{#1}{#2}{#3}}%
}
\makeatother
